% Source: https://artofproblemsolving.com/community/c3h2478663_egyptian_fractions
% Stable identity: aops-topic-2478663-q01
\textbf{A general method.} For positive integers $p,q$, multiplying $\dfrac1x+\dfrac1y=\dfrac pq$ by $qxy$ gives $qy+qx=pxy$, i.e. $p^2xy-pqx-pqy=0$, and adding $q^2$ to both sides,
\[(px-q)(py-q)=q^2.\]
So the solutions come from the factor pairs of $q^2$.

\textbf{Here} $p=3$, $q=10$:
\[(3x-10)(3y-10)=100.\]
Since $\dfrac1x<\dfrac3{10}$, we have $3x>10$, and similarly $3y>10$, so both factors are positive divisors of $100$. Also $3x-10\equiv-10\equiv2\pmod3$, and the same holds for $3y-10$. The divisors of $100$ are $1,2,4,5,10,20,25,50,100$; those congruent to $2$ modulo $3$ are $2,5,20,50$. The factor pairs of $100$ made of such divisors, with the smaller factor first, are
\[2\cdot50\qquad\text{and}\qquad5\cdot20.\]
They give $(x,y)=(4,20)$ and $(x,y)=(5,10)$. Check: $\dfrac14+\dfrac1{20}=\dfrac{6}{20}=\dfrac3{10}$ and $\dfrac15+\dfrac1{10}=\dfrac3{10}$.

The pairs are $\boxed{(4,20)\text{ and }(5,10)}$.
